\label{Section:Rust ported}
Some modification have been made in order to support the Rust language.

\subsection{Limitation of Lazart for Rust}

Although Lazart works on LLVM and rustc is capable of producing LLVM, Some limitation was encountered, when trying to analyse Rust programmers with Lazart.

\textbf{LLVM version.} The LLVM version of Lazart was updated from LLVM 9 to LLVM 14, as the Rust tool chain matching LLVM 9 is too old for cargo to fetch crates from registry.

\textbf{Test inversion extension.} rustc makes an extensive use of LLVM's \texttt{select} and \texttt{switch} instructions compared to C compilers (here \texttt{clang}).
Rust is capable to generate \texttt{switch} and \texttt{select} from simple branching (example: condition branch), even if the \texttt{switch} instructions is often used for pattern matching.
The test inversion model has been extended to support \texttt{switch} and \texttt{select} LLVM instructions in order to be able to fault some control flow constructs (to match C example in FISSC) or some language constructs (e.g. pattern matching). 

\subsection{The Change Pointer fault model extension}
\label{subsection:Change Pointer fault model}
A new form of attack as been added.
Since Rust borrow checker guarantees safety at compiled time. Nothing guarantees those safety to protected at run-time are they are purely static.
Since developers rely on the borrow checker for their program and to enable optimisation. One way to break the borrow checker guarantees is to consider a fault model of \textit{change pointer} making possible to redirect a mutable reference to another.
This fault is made possible with the help of the \texttt{unsafe} keyword.
At lower representation levels, the fault can map a skip on an address computation, or a fault on a register or the stack.

\textbf{Change Pointer fault model} The Change Pointer (CP) fault model, implemented in Lazart for Rust and C, redirects a reference so that it point to another object.
Fault yielding a dangling or wild pointer are out of scope due to them raising a trap and become easily detectable.
A redirection between two legitimate references is silent and is exactly what defeats a borrow checker invariant.

\textbf{Example: privilege escalation by reference swap.} Listing \ref{lst:cp-borrow} shows a typical use-case example of the model in Lazart.The target program uses the function \texttt{connect} to legitimately set the session with a specified duration.The two \textit{Change Pointer} injection points (using the macro \texttt{\_LZ\_\_change\_ptr\_R!(good,\allowbreak{} bad,\allowbreak{} id,\allowbreak{} $\tau$)}) modify the nominal parameters \texttt{connected} and \texttt{duration} to \texttt{admin} and \texttt{sold} respectively.

\begin{lstlisting}[language=Rust, caption={Privilege escalation via change-pointer faults.}, label={lst:cp-borrow}]
fn connect(session: &mut bool, amount: &mut i32) {
    *session = true;   // legitimate, non-privileged write
    *amount  = 100;
}

fn main() {
    let (mut admin, mut connected) = (false, false);
    let (mut sold,  mut duration)  = (0, 0);

    // Pointer change on function call.
    connect(_LZ__change_ptr_R!(&mut connected, &mut admin, "is_admin",  bool), 
            _LZ__change_ptr_R!(&mut duration,  &mut sold,  "have_sold", i32));
    _LZ__ORACLE!(admin & (sold == 100)); // always false without faults.
}
\end{lstlisting}

Lazart finds an attack combining a fault on each injection point to force a mutable aliasing that is forbidden by Rust, leading to privilege escalation.
Although expected, these results show that the memory safety \textit{by construction} does not imply \textit{fault injection} resistance.

\textbf{Note:} the implementation of such faults is not present in this paper (see Disclaimer at the end) and is discussed in the FDTC paper (under submission) \cite{fdct2026}. 